\documentclass[12pt]{article}
\usepackage[utf8]{inputenc}
\usepackage[T1]{fontenc}
\usepackage{hyperref}
\title{Jo Language Resources for K-12 Curriculum Development}
\author{Compiled by Claude Sonnet 5.5 (claude-sonnet-5-5) for Andr\'as Kornai}
\date{}
\begin{document}
\maketitle

\section{Introduction}
Jo (ISO 639-3 wsi, Glottocode wusi1237) is an Austronesian language of Vanuatu. The same code covers the language usually called Wusi, an Oceanic language of the west coast of Espiritu Santo. One catalogue gives the alternate names Wusi-Kerepua and Kula and about 350 speakers. Global Recordings Network lists two Jo varieties under this code: Jo: Elia and Jo: Kula Big Bay Bush. It is not known to be taught in schools.

\section{Primary Educational Resources}
No primers, textbooks or school programmes were found.

\section{Linguistic and Reference Materials}
\begin{itemize}
\item Wikipedia, Wusi language. \url{https://en.wikipedia.org/wiki/Wusi_language}
\item Verbix Endangered Languages page for Wusi. \url{https://docs.verbix.com/EndangeredLanguages/Wusi}
\item Endangered Languages catalogue entries. \url{https://endangeredlanguages.com/lang/2518} and \url{https://www.endangeredlanguages.com/lang/2511}
\item Endangered Languages bibliography page for this language. \url{https://endangeredlanguages.com/lang/3321/bibliography?hl=ru}
\item An article in Language Documentation \& Description appeared in the search results. I did not verify that it concerns this language. \url{https://lddjournal.org/article/1183/galley/2428/download}
\end{itemize}

\section{Religious and Community Resources}
No scripture translations or community organisations were found.

\section{Digital and Audio Resources}
Global Recordings Network pages for the two Jo varieties: Elia \url{https://glorecnet-website.k.globalrecordings.net/en/language/29971} and Kula Big Bay Bush \url{https://glorecnet-website.k.globalrecordings.net/en/language/29972}. I did not open them, so I cannot say which recordings they contain.

\section{Curriculum Development Context}
Little published material exists. Vanuatu schooling is in Bislama, English and French, so a Jo curriculum would need to be built with the community from scratch, starting with documentation and a practical orthography. Audio from Global Recordings Network may serve as a starting corpus.
\end{document}